\documentclass[12pt,openany]{book}
\usepackage{amsmath}
\usepackage{amssymb}
\usepackage{amsthm}
\usepackage{mathtools}
\usepackage[hidelinks]{hyperref}
\usepackage{booktabs}
\usepackage{geometry}
\usepackage{fancyhdr}
\usepackage{setspace}
\usepackage[expansion=false]{microtype}
\usepackage{xcolor}
\usepackage{enumitem}
\usepackage{titlesec}
\usepackage{tocloft}
\usepackage{array}
\usepackage{tabularx}
\usepackage{float}
\usepackage{longtable}
\usepackage{ragged2e}
\usepackage{thmtools}
\usepackage{chngcntr}
\counterwithout{table}{chapter}
\counterwithout{figure}{chapter}
\makeatletter
\renewcommand{\theHtable}{\number\value{table}}
\renewcommand{\theHfigure}{\number\value{figure}}
\makeatother
\usepackage{sectionbreak}
\usepackage{algorithmicx}
\usepackage{algpseudocode}
\usepackage{algorithm}

% Guard against overfull hboxes in bibliography entries
\usepackage{etoolbox}
\apptocmd{\thebibliography}{\raggedright}{}{}

% ============================================================
% CUSTOM MACROS
% ============================================================
\newcommand{\R}{\mathbb{R}}
\newcommand{\E}{\mathbb{E}}
\newcommand{\Prob}{\mathbb{P}}
\newcommand{\N}{\mathbb{N}}

\newcommand{\bfw}{\mathbf{w}}
\newcommand{\bfx}{\mathbf{x}}
\newcommand{\bfz}{\mathbf{z}}
\newcommand{\bfy}{\mathbf{y}}
\newcommand{\bfmu}{\boldsymbol{\mu}}
\newcommand{\bfSigma}{\boldsymbol{\Sigma}}
\newcommand{\bfQ}{\mathbf{Q}}
\newcommand{\bfR}{\mathbf{R}}
\newcommand{\bfF}{\mathbf{F}}
\newcommand{\bfH}{\mathbf{H}}
\newcommand{\bfK}{\mathbf{K}}
\newcommand{\bfP}{\mathbf{P}}
\newcommand{\bfI}{\mathbf{I}}
\newcommand{\bfA}{\mathbf{A}}
\newcommand{\bfe}{\mathbf{e}}
\newcommand{\bfv}{\mathbf{v}}
\newcommand{\bfu}{\mathbf{u}}
\newcommand{\bfGamma}{\boldsymbol{\Gamma}}
\newcommand{\bfOmega}{\boldsymbol{\Omega}}
\newcommand{\bfxi}{\boldsymbol{\xi}}
\newcommand{\bfone}{\mathbf{1}}

\newcommand{\calE}{\mathcal{E}}
\newcommand{\calF}{\mathcal{F}}
\newcommand{\calG}{\mathcal{G}}
\newcommand{\calP}{\mathcal{P}}
\newcommand{\calU}{\mathcal{U}}
\newcommand{\calM}{\mathcal{M}}
\newcommand{\calT}{\mathcal{T}}
\newcommand{\calR}{\mathcal{R}}
\newcommand{\calS}{\mathcal{S}}
\newcommand{\calI}{\mathcal{I}}
\newcommand{\calX}{\mathbf{X}}
\newcommand{\calA}{\mathcal{A}}
\newcommand{\calO}{\mathcal{O}}
\newcommand{\calC}{\mathcal{C}}
\newcommand{\calL}{\mathcal{L}}

\newcommand{\MaxDD}{\mathrm{MaxDD}}
\newcommand{\DD}{\mathrm{DD}}

\DeclareMathOperator{\CVaR}{CVaR}
\DeclareMathOperator{\VaR}{VaR}
\DeclareMathOperator{\Var}{Var}
\DeclareMathOperator{\Cov}{Cov}
\DeclareMathOperator*{\argmax}{arg\,max}
\DeclareMathOperator*{\argmin}{arg\,min}
\DeclareMathOperator{\tr}{tr}
\DeclareMathOperator{\rank}{rank}

% ============================================================
% GEOMETRY AND LAYOUT
% ============================================================
\geometry{margin=1in, includefoot}
\microtypesetup{expansion=false, final}
\singlespacing
\setlength{\parindent}{0.5in}
\setlength{\parskip}{0.5ex plus 0.2ex minus 0.1ex}
\setlength{\emergencystretch}{1.5em}
\tolerance=9999
\hbadness=10000
\vbadness=10000
\hyphenpenalty=10000
\exhyphenpenalty=10000
\raggedbottom
\sloppy

\setlength{\abovedisplayskip}{6pt plus 2pt minus 3pt}
\setlength{\belowdisplayskip}{6pt plus 2pt minus 3pt}
\setlength{\abovedisplayshortskip}{3pt plus 1pt minus 1pt}
\setlength{\belowdisplayshortskip}{3pt plus 1pt minus 1pt}
\numberwithin{equation}{chapter}
\allowdisplaybreaks[3]

% ============================================================
% THEOREM ENVIRONMENTS
% ============================================================
\theoremstyle{plain}
\newtheorem{theorem}{Theorem}[chapter]
\newtheorem{lemma}[theorem]{Lemma}
\newtheorem{corollary}[theorem]{Corollary}
\newtheorem{proposition}[theorem]{Proposition}
\newtheorem{conjecture}[theorem]{Conjecture}
\theoremstyle{definition}
\newtheorem{definition}[theorem]{Definition}
\newtheorem{axiom}[theorem]{Axiom}
\newtheorem{assumption}[theorem]{Assumption}
\theoremstyle{remark}
\newtheorem{remark}[theorem]{Remark}

\providecommand{\theHtheorem}{}
\renewcommand{\theHtheorem}{\thechapter-\arabic{theorem}}
\providecommand{\theHlemma}{}
\renewcommand{\theHlemma}{\thechapter-\arabic{theorem}}
\providecommand{\theHcorollary}{}
\renewcommand{\theHcorollary}{\thechapter-\arabic{theorem}}
\providecommand{\theHproposition}{}
\renewcommand{\theHproposition}{\thechapter-\arabic{theorem}}
\providecommand{\theHdefinition}{}
\renewcommand{\theHdefinition}{\thechapter-\arabic{theorem}}
\providecommand{\theHaxiom}{}
\renewcommand{\theHaxiom}{\thechapter-\arabic{theorem}}
\providecommand{\theHassumption}{}
\renewcommand{\theHassumption}{\thechapter-\arabic{theorem}}
\providecommand{\theHremark}{}
\renewcommand{\theHremark}{\thechapter-\arabic{theorem}}
\providecommand{\theHconjecture}{}
\renewcommand{\theHconjecture}{\thechapter-\arabic{theorem}}

\makeatletter
\providecommand*{\toclevel@conjecture}{0}
\makeatother

\newenvironment{abstract}{%
  \chapter*{Abstract}
  \begin{quote}\small
}{%
  \end{quote}
}
\newenvironment{keywords}{%
  \paragraph*{Keywords:}
}{}

% ============================================================
% SECTION FORMATTING
% ============================================================
\titleformat{\chapter}[display]
    {\normalfont\Huge\bfseries}{\chaptertitlename\ \thechapter}{20pt}{\Huge}
\titlespacing*{\chapter}{0pt}{50pt}{40pt}
\titleformat{\section}{\normalfont\Large\bfseries}{\thesection}{1em}{}
\titlespacing*{\section}{0pt}{3.5ex plus 1ex minus 0.2ex}{2.3ex plus 0.2ex}
\titleformat{\subsection}{\normalfont\large\bfseries}{\thesubsection}{1em}{}
\titlespacing*{\subsection}{0pt}{3.25ex plus 1ex minus 0.2ex}{1.5ex plus 0.2ex}
\titleformat{\subsubsection}{\normalfont\normalsize\bfseries}{\thesubsubsection}{1em}{}
\titlespacing*{\subsubsection}{0pt}{3ex plus 0.5ex minus 0.2ex}{1.5ex plus 0.2ex}

% ============================================================
% TABLE OF CONTENTS
% ============================================================
\renewcommand{\cftchapleader}{\cftdotfill{\cftdotsep}}
\setlength{\cftchapindent}{0em}
\setlength{\cftsecindent}{1.5em}

% ============================================================
% HEADER AND FOOTER
% ============================================================
\pagestyle{fancy}
\fancyhf{}
\fancyhead[LE]{\thepage}
\fancyhead[RO]{\thepage}
\fancyhead[RE]{\nouppercase{\leftmark}}
\fancyhead[LO]{\nouppercase{\rightmark}}
\renewcommand{\headrulewidth}{0.4pt}
\renewcommand{\footrulewidth}{0pt}
\setlength{\headheight}{27.2pt}

% ============================================================
% MISCELLANEOUS
% ============================================================
\setlist[enumerate]{leftmargin=*, partopsep=0.1ex}
\clubpenalty=10000
\widowpenalty=10000
\displaywidowpenalty=10000
\renewcommand{\arraystretch}{1.2}
\setlength{\tabcolsep}{4pt}

\pdfstringdefDisableCommands{%
  \def\mathcal#1{#1}%
  \def\bfz{z}%
  \def\calE{E}%
  \def\calF{F}%
  \def\calG{G}%
  \def\calP{P}%
  \def\calU{U}%
  \def\calM{M}%
  \def\calT{T}%
  \def\calR{R}%
  \def\calS{S}%
  \def\calI{I}%
}

\title{Team-Mates Whitepaper:\\
Multi-Agent Quantitative-Finance Platform\\[0.6em]
{\large Alpaca Trading Hackathon 2026}}
\author{Hadrian Hu}
\date{\today}

\begin{document}

\maketitle

\tableofcontents
\listoftheorems[ignoreall,show={theorem,lemma,proposition,corollary,definition,axiom,assumption,conjecture}]

% ==============================================================
\begin{abstract}
This whitepaper presents the \emph{Multi-Agent Quantitative-Finance
Platform}, a finite-capital investment architecture designed for
the Alpaca Trading Hackathon 2026.  The platform integrates
heterogeneous specialised agents across four financial domains ---
economic, financial, fiscal, and investment --- into a unified
capital allocation engine governed by a seven-dimensional
state-of-charge vector and a Kalman-filter prediction--correction
cycle.  The central thesis is that optimal agentic trading is
structurally isomorphic to minimum-variance unbiased estimation
(MVUE) of a latent market state given partial, noisy, multi-domain
observations.  The finite initial capital of \$100{,}000 is not a
handicap but a precision condition: it forces the system to operate
in the well-calibrated, survival-first regime that preserves the
principal necessary for long-run compounding.  The platform is
designed to be fully auditable, reproducible, and extensible.
\end{abstract}

\begin{keywords}
multi-agent systems; Kalman filter; state-space model; capital
allocation; finite initial capital; state-of-charge; seven-state
architecture; Bayesian reputation; ensemble signal; innovation
covariance; circuit-breaker; risk-adjusted objective; portfolio
diversification; Alpaca hackathon
\end{keywords}

% ==============================================================
\chapter*{Executive Summary}
\addcontentsline{toc}{chapter}{Executive Summary}
% ==============================================================

\noindent
The platform answers one research question:

\begin{quote}
\emph{Given all information legitimately available at time $t$,
what market regime is most probable, what risks currently dominate,
how uncertain is that assessment, and what portfolio allocation
provides the best risk-adjusted response under a \$100{,}000 finite
capital constraint?}
\end{quote}

\noindent
The answer is formalised as:

\begin{equation}
  \begin{aligned}
    \text{Optimal Agentic Trading}
    \;=\;\;&\text{MVUE of Latent Market State Given}\\
    &\text{Partial, Noisy, Multi-Domain Observations}
  \end{aligned}
\end{equation}

\noindent\textbf{Priority ordering.}

\begin{equation}
  \text{Survival}
  \;\succ\; \text{Preservation}
  \;\succ\; \text{Intelligent Allocation}
  \;\succ\; \text{Controlled Risk}
  \;\succ\; \text{Compounding}
\end{equation}

\noindent\textbf{Architecture at a glance.}

\begin{center}
\renewcommand{\arraystretch}{1.35}
\begin{tabularx}{\textwidth}{@{} l X @{}}
\toprule
\textbf{Layer} & \textbf{Function} \\
\midrule
Data ingestion & Historical and real-time market, macro, and news feeds. \\
Specialised agents & Independent analysis across four domains; emit $(s_i, c_i, u_i, d_i)$. \\
State estimation & Seven-dimensional SoC vector $\mathbf{X}_t \in [0,1]^7$. \\
Kalman filter & Prediction--correction cycle; Kalman gain $K_t$ as deployment weight. \\
State gating & Product gate $G(\mathbf{X}_t)$; circuit-breaker on critical threshold breach. \\
Capital allocator & $\pi^* = K_t \cdot C_{\mathrm{available}} \cdot G(\mathbf{X}_t)$. \\
Risk gate & CVaR, MaxDD, concentration constraints. \\
Execution & Alpaca API paper-trading orders. \\
Performance & Composite objective $J$ vs.\ five baselines. \\
Reputation update & Bayesian $w_i$ adaptation after each outcome. \\
\bottomrule
\end{tabularx}
\end{center}

% ==============================================================
\chapter{Platform Architecture and Design Philosophy}
\label{ch:architecture}
% ==============================================================

\section{Core Principle}

\noindent
Success is not defined as maximising predicted profit.  It is
defined as finding decisions that remain defensible when models
disagree, markets change regime, forecasts fail, correlations
break, and uncertainty becomes materially greater.

\section{The Four Financial Domains}

The platform spans four interacting financial domains, treated as
four views of one system rather than four independent applications:

\begin{center}
\renewcommand{\arraystretch}{1.35}
\begin{tabularx}{\textwidth}{@{} l l X @{}}
\toprule
\textbf{Domain} & \textbf{Symbol} & \textbf{Scope} \\
\midrule
Economic & $\calE$ & GDP, inflation, employment, credit, fiscal impulse. \\
Financial & $\calF$ & Asset prices, volatility surface, yield curve, spreads. \\
Fiscal & $\calG$ & Government expenditure, tax regime, deficit, debt. \\
Investment & $\calI$ & Fund flows, positioning, sentiment, options markets. \\
\bottomrule
\end{tabularx}
\end{center}

\section{Initial Capital Partition}

\begin{definition}[Capital Partition]
\label{def:partition}
The initial capital $W_0 = \$100{,}000$ is partitioned into six
disjoint buckets $\{b_k\}_{k=1}^{6}$ satisfying:
\begin{enumerate}[label=\textbf{(P\arabic*)}]
  \item $\sum_{k=1}^{6} b_k = W_0$ (exhaustiveness).
  \item $b_k \geq 0$ for all $k$ (non-negativity).
  \item At least one $b_k > 0$ (non-triviality).
\end{enumerate}
\end{definition}

\begin{table}[H]
\centering
\caption{Canonical six-bucket partition of the \$100{,}000 initial capital.}
\label{tab:buckets}
\small
\begin{tabularx}{\linewidth}{@{} l r r X @{}}
\toprule
\textbf{Bucket} & \textbf{Amount (\$)} & \textbf{Weight}
  & \textbf{Purpose} \\
\midrule
Core / Stable               & 35{,}000 & 0.35
  & Capital preservation; low-volatility instruments. \\
Diversified Growth          & 20{,}000 & 0.20
  & Broad-market multi-asset growth exposure. \\
Opportunistic               & 15{,}000 & 0.15
  & Tactical short-horizon strategies. \\
Alternative / Higher-Vol    & 10{,}000 & 0.10
  & Asymmetric or higher-risk instruments. \\
Cash / Liquidity Reserve    & 10{,}000 & 0.10
  & Mandatory floor; not deployed without trigger. \\
Dynamic Agent Allocation    & 10{,}000 & 0.10
  & Autonomously managed; agents justify movements. \\
\midrule
\textbf{Total}              & \textbf{100{,}000} & \textbf{1.00} & \\
\bottomrule
\end{tabularx}
\end{table}

\begin{axiom}[Mandatory Liquidity Floor]
\label{ax:liquidity-floor}
At all times $t$, the cash reserve satisfies:
\begin{equation}
  W_t^{(5)} \;\geq\; \ell_{\min} \;:=\; 0.05 \cdot W_0 \;=\; \$5{,}000.
  \label{eq:liquidity-floor}
\end{equation}
No signal overrides this constraint without explicit risk-gate
approval and a documented audit entry.
\end{axiom}

% ==============================================================
\chapter{Agent Output Contract and Ensemble Aggregation}
\label{ch:agents}
% ==============================================================

\section{Agent Output Tuple}

\begin{definition}[Agent Output Tuple]
\label{def:agent-tuple}
At each time $t$, specialised agent $i \in \{1,\ldots,N\}$ emits:
\begin{equation}
  \calA_i(t)
  \;=\;
  \bigl(\,s_i,\;\; c_i,\;\; u_i,\;\; d_i,\;\;
         p_i^{+},\;\; p_i^{-},\;\; \Delta t_i,\;\; r_i\,\bigr)
  \label{eq:agent-tuple}
\end{equation}
with domains: $s_i \in [-1,1]$, $c_i \in (0,1]$,
$u_i \in [0,1]$, $d_i \in [0,1]$, $p_i^+,p_i^- \geq 0$ with
$p_i^+ + p_i^- \leq 1$, $\Delta t_i > 0$, $r_i \geq 0$.
\end{definition}

\section{Ensemble Aggregate Signal}

\begin{definition}[Ensemble Aggregate Signal]
\label{def:ensemble-signal}
The confidence-weighted, uncertainty-discounted, doubt-dampened
aggregate signal is:
\begin{equation}
  S_t
  \;=\;
  \frac{\displaystyle\sum_{i=1}^{N}
        w_i\,s_i\,c_i\,(1-u_i)\,(1-d_i)}
       {\displaystyle\sum_{i=1}^{N} w_i},
  \label{eq:ensemble-signal}
\end{equation}
where $w_i > 0$ are Bayesian reputation weights.
\end{definition}

\begin{theorem}[Boundedness of $S_t$]
\label{thm:signal-bounded}
Under Definition~\ref{def:agent-tuple} and
Definition~\ref{def:ensemble-signal}, $S_t \in [-1,1]$.
\end{theorem}

\begin{proof}
Each dampened signal $\phi_i = s_i c_i(1-u_i)(1-d_i)$ satisfies
$|\phi_i| \leq |s_i| \leq 1$ since $c_i, (1-u_i), (1-d_i) \in [0,1]$.
Therefore $S_t$ is a convex combination of values in $[-1,1]$ and
thus $S_t \in [-1,1]$.
\qed
\end{proof}

\begin{definition}[Agent Disagreement]
\label{def:disagreement}
The weighted mean absolute deviation of agent signals from the
ensemble mean is:
\begin{equation}
  D_t \;=\;
  \frac{\displaystyle\sum_{i=1}^{N} w_i\,|s_i - S_t|}
       {\displaystyle\sum_{i=1}^{N} w_i}.
  \label{eq:disagreement}
\end{equation}
$D_t \geq 0$ always; $D_t = 0$ iff all agents agree.
\end{definition}

\begin{remark}[Semantic Significance of Disagreement]
$S_t = 0$ can arise from genuine neutrality ($D_t \approx 0$) or
from violent disagreement ($D_t$ large).  Only $D_t$ distinguishes
these two fundamentally different risk states.  High $D_t$ inflates
the innovation covariance, reducing the Kalman gain and limiting
capital deployment regardless of the directional signal.
\end{remark}

% ==============================================================
\chapter{The Kalman Filter Investment Architecture}
\label{ch:kalman}
% ==============================================================

\section{Structural Isomorphism}

The central theoretical result is the following bijection between
the Kalman prediction--correction cycle and the investment
allocation cycle:

\begin{table}[H]
\centering
\caption{Kalman filter -- investment allocation isomorphism.}
\label{tab:isomorphism}
\small
\begin{tabularx}{\textwidth}{@{} p{5.5cm} X @{}}
\toprule
\textbf{Kalman Filter Concept} & \textbf{Investment Architecture Concept} \\
\midrule
Hidden state $\hat{\mathbf{X}}_{t|t-1}$ & Prior expected return and risk model \\
Prediction covariance $\mathbf{P}_{t|t-1}$ & Pre-signal portfolio variance \\
Observation $\mathbf{z}_t$ & Agent signals $(s_i, c_i, u_i, d_i)$ \\
Innovation $\tilde{\mathbf{y}}_t$ & Aggregate signal minus prior forecast \\
Kalman gain $\mathbf{K}_t$ & Confidence-weighted capital deployment fraction \\
Posterior estimate $\hat{\mathbf{X}}_{t|t}$ & Updated portfolio weights $\mathbf{w}_{t|t}$ \\
Posterior covariance $\mathbf{P}_{t|t}$ & Post-signal residual portfolio risk \\
Process noise $\mathbf{Q}$ & Irreducible market uncertainty (aleatoric) \\
Measurement noise $\mathbf{R}_t$ & Agent estimation error and calibration doubt \\
\bottomrule
\end{tabularx}
\end{table}

\section{The Kalman Prediction--Correction Cycle}

\textbf{Prediction step} (time update):
\begin{align}
  \hat{\mathbf{X}}_{t|t-1} &= \mathbf{F}\,\hat{\mathbf{X}}_{t-1|t-1}, \\
  \mathbf{P}_{t|t-1} &= \mathbf{F}\,\mathbf{P}_{t-1|t-1}\,\mathbf{F}^\top + \mathbf{Q}.
\end{align}

\textbf{Correction step} (measurement update):
\begin{align}
  \tilde{\mathbf{y}}_t &= \mathbf{z}_t - \mathbf{H}\,\hat{\mathbf{X}}_{t|t-1}, \\
  \mathbf{K}_t &= \mathbf{P}_{t|t-1}\,\mathbf{H}^\top
    \!\left(\mathbf{H}\,\mathbf{P}_{t|t-1}\,\mathbf{H}^\top + \mathbf{R}_t\right)^{-1}, \\
  \hat{\mathbf{X}}_{t|t} &= \hat{\mathbf{X}}_{t|t-1} + \mathbf{K}_t\,\tilde{\mathbf{y}}_t, \\
  \mathbf{P}_{t|t} &= \left(\mathbf{I} - \mathbf{K}_t\,\mathbf{H}\right)\mathbf{P}_{t|t-1}.
\end{align}

\section{The Investment Kalman Gain}

The ensemble effective confidence is:
\begin{equation}
  \bar{c}_t \;=\;
  \frac{\displaystyle\sum_{i=1}^{N} w_i\,c_i\,(1-u_i)\,(1-d_i)}
       {\displaystyle\sum_{i=1}^{N} w_i}
  \;\in\; [0,1].
  \label{eq:effective-confidence}
\end{equation}

The effective measurement noise, with disagreement inflation, is:
\begin{equation}
  R_t \;=\;
  \sigma_{\mathrm{base}}^2\!\cdot\!\frac{1-\bar{c}_t}{\bar{c}_t}
  \;+\;
  D_t^2\,\sigma_{\mathrm{base}}^2.
  \label{eq:effective-noise}
\end{equation}

The scalar investment Kalman gain is:
\begin{equation}
  K_t^{(\mathrm{inv})} \;=\;
  \frac{P_{t|t-1}}{P_{t|t-1} + R_t}
  \;\in\; [0,1).
  \label{eq:investment-gain}
\end{equation}

\begin{proposition}[Properties of $K_t^{(\mathrm{inv})}$]
\label{prop:gain-properties}
The investment Kalman gain satisfies:
\begin{enumerate}
  \item $K_t^{(\mathrm{inv})} \in [0,1)$ for all $t$ (since $R_t \geq 0$
        and $P_{t|t-1} > 0$).
  \item $K_t^{(\mathrm{inv})} \to 1$ as $\bar{c}_t \to 1$, $D_t \to 0$
        (perfect ensemble).
  \item $K_t^{(\mathrm{inv})} \to 0$ as $\bar{c}_t \to 0$ or $D_t \to \infty$
        (no usable signal).
  \item $K_t^{(\mathrm{inv})}$ is strictly decreasing in $R_t$ and strictly
        increasing in $P_{t|t-1}$.
\end{enumerate}
\end{proposition}

% ==============================================================
\chapter{The Seven-State Investment Architecture}
\label{ch:states}
% ==============================================================

\section{The Seven-Dimensional Hidden State}

\begin{definition}[Seven-Dimensional Investment State]
\label{def:hidden-state}
The hidden investment state vector at time $t$ is:
\begin{equation}
  \mathbf{X}_t \;=\;
  \bigl(\calE_t,\;\calF_t,\;\calG_t,\;\calP_t,\;
         \calU_t,\;\calM_t,\;\calT_t\bigr)^\top
  \;\in\; [0,1]^7,
  \label{eq:hidden-state}
\end{equation}
where each component $\mathcal{S}_t^{(d)} \in [0,1]$ is the
State-of-Charge (SoC) of the corresponding investment dimension.
\end{definition}

\begin{table}[H]
\centering
\caption{The seven investment state dimensions.}
\label{tab:seven-states}
\small
\begin{tabularx}{\textwidth}{@{} l l X X @{}}
\toprule
\textbf{Symbol} & \textbf{Name} & \textbf{Full Charge ($\approx 1$)}
  & \textbf{Discharged ($\approx 0$)} \\
\midrule
$\calE_t$ & Economic    & Expansion, strong growth & Recession, contraction \\
$\calF_t$ & Financial   & Low stress, tight spreads & Crisis, illiquid \\
$\calG_t$ & Fiscal      & Expansionary fiscal policy & Fiscal cliff, austerity \\
$\calP_t$ & Portfolio   & Low drawdown, ample cash & Near drawdown limit \\
$\calU_t$ & Fundamental & Reasonable valuations & Bubble or distress \\
$\calM_t$ & Market      & Deep, orderly markets & Thin, erratic \\
$\calT_t$ & Sector/Tech & Strong secular tailwinds & Structural headwinds \\
\bottomrule
\end{tabularx}
\end{table}

\section{Charge Dynamics and Discharge Asymmetry}

\begin{definition}[State Charge Dynamics]
\label{def:charge-dynamics}
For each dimension $d$ and each time step $t$:
\begin{equation}
  \mathcal{S}_{t+1}^{(d)} \;=\; \Pi\!\left(
  \begin{cases}
    \mathcal{S}_t^{(d)} + \kappa_d^{+}(1 - \mathcal{S}_t^{(d)})
      + w_t^{(d)} & E_t^{+} \\
    \mathcal{S}_t^{(d)}(1 - \kappa_d^{-}) + w_t^{(d)} & E_t^{-} \\
    \mathcal{S}_t^{(d)} + w_t^{(d)} & E_t^{0}
  \end{cases}
  \right)
  \label{eq:charge-dynamics}
\end{equation}
where $\Pi(x) = \max(0,\min(1,x))$ is the unit-interval
projection, $\kappa_d^{+}$ and $\kappa_d^{-}$ are the charge and
discharge rates, and $w_t^{(d)} \sim \mathcal{N}(0,q_d^2)$ is
process noise.
\end{definition}

\begin{theorem}[Discharge Dominance]
\label{thm:discharge-dominance}
Under symmetric shocks ($p = \tfrac{1}{2}$) and discharge
asymmetry $\kappa_d^{-} > \kappa_d^{+}$, the long-run expected
state charge satisfies:
\begin{equation}
  \E\!\left[\mathcal{S}_\infty^{(d)}\right]
  \;=\;
  \frac{\kappa_d^{+}}{\kappa_d^{+} + \kappa_d^{-}}
  \;<\; \frac{1}{2}.
  \label{eq:discharge-dominance}
\end{equation}
Financial systems discharge faster than they recover.
\end{theorem}

\begin{table}[H]
\centering
\caption{Charge and discharge rates with asymmetry ratios.}
\label{tab:charge-rates}
\small
\begin{tabularx}{\textwidth}{@{} l l l l l @{}}
\toprule
\textbf{State} & \textbf{$\kappa_d^{+}$} & \textbf{$\kappa_d^{-}$}
  & \textbf{Asymmetry} & \textbf{$\E[\mathcal{S}_\infty^{(d)}]$} \\
\midrule
Economic    & 0.03/month & 0.08/month & $2.7\times$ & 0.273 \\
Financial   & 0.15/day   & 0.60/day   & $4.0\times$ & 0.200 \\
Fiscal      & 0.01/month & 0.02/month & $2.0\times$ & 0.333 \\
Portfolio   & 0.05/day   & 0.20/day   & $4.0\times$ & 0.200 \\
Fundamental & 0.02/month & 0.04/month & $2.0\times$ & 0.333 \\
Market      & 0.50/hour  & 2.0/hour   & $4.0\times$ & 0.200 \\
Sector/Tech & 0.005/month & 0.010/month & $2.0\times$ & 0.333 \\
\bottomrule
\end{tabularx}
\end{table}

% ==============================================================
\chapter{The State Gating Architecture}
\label{ch:gating}
% ==============================================================

\section{Critical and Full-Charge Thresholds}

\begin{definition}[Critical and Full-Charge Thresholds]
\label{def:thresholds}
For each $d \in \{1,\ldots,7\}$, the critical threshold
$\mathcal{S}_{\min}^{(d)}$ and full-charge threshold
$\mathcal{S}_{\mathrm{full}}^{(d)}$ satisfy:
\begin{equation}
  0 < \mathcal{S}_{\min}^{(d)} < \mathcal{S}_{\mathrm{full}}^{(d)} \leq 1.
\end{equation}
\end{definition}

\begin{table}[H]
\centering
\caption{Canonical threshold values for the seven state dimensions.}
\label{tab:thresholds}
\small
\begin{tabularx}{\textwidth}{@{} l l l l l @{}}
\toprule
\textbf{Dim.} & \textbf{State} & $\mathcal{S}_{\min}^{(d)}$
  & $\mathcal{S}_{\mathrm{full}}^{(d)}$ & \textbf{Consequence of breach} \\
\midrule
1 & $\calE_t$ Economic    & 0.15 & 0.70 & Capital preservation mode \\
2 & $\calF_t$ Financial   & 0.20 & 0.80 & All new positions blocked \\
3 & $\calG_t$ Fiscal      & 0.10 & 0.65 & Long-horizon allocations frozen \\
4 & $\calP_t$ Portfolio   & 0.20 & 0.70 & Hard stop; drawdown limit \\
5 & $\calU_t$ Fundamental & 0.15 & 0.75 & Reduce equity; seek confirmation \\
6 & $\calM_t$ Market      & 0.20 & 0.80 & No new orders; monitor only \\
7 & $\calT_t$ Sector/Tech & 0.15 & 0.75 & Exit disrupted sectors \\
\bottomrule
\end{tabularx}
\end{table}

\section{The Gating Function}

\begin{definition}[State Gating Function]
\label{def:gating-function}
For each dimension $d$, the piecewise-linear gating function is:
\begin{equation}
  g_d(\mathcal{S}) \;=\; \begin{cases}
    0
      & \mathcal{S} < \mathcal{S}_{\min}^{(d)}, \\[4pt]
    \displaystyle\frac{\mathcal{S} - \mathcal{S}_{\min}^{(d)}}
                      {\mathcal{S}_{\mathrm{full}}^{(d)} - \mathcal{S}_{\min}^{(d)}}
      & \mathcal{S} \in \bigl[\mathcal{S}_{\min}^{(d)},\,
                               \mathcal{S}_{\mathrm{full}}^{(d)}\bigr), \\[4pt]
    1
      & \mathcal{S} \geq \mathcal{S}_{\mathrm{full}}^{(d)}.
  \end{cases}
  \label{eq:gating-function}
\end{equation}
\end{definition}

\begin{proposition}[Properties of $g_d$]
\label{prop:gating-properties}
$g_d : [0,1] \to [0,1]$ is non-decreasing, Lipschitz continuous,
and satisfies the circuit-breaker biconditional:
\begin{equation}
  g_d(\mathcal{S}) = 0
  \;\Longleftrightarrow\;
  \mathcal{S} \leq \mathcal{S}_{\min}^{(d)}.
\end{equation}
\end{proposition}

\section{The Product Gating Function}

\begin{definition}[Product Gating Function]
\label{def:product-gating}
The seven-dimensional product gate is:
\begin{equation}
  G(\mathbf{X}_t) \;=\; \prod_{d=1}^{7} g_d\!\left(\mathcal{S}_t^{(d)}\right)
  \;\in\; [0,1].
  \label{eq:product-gating}
\end{equation}
\end{definition}

\begin{theorem}[Product Gating Properties]
\label{thm:product-gating}
$G(\mathbf{X}_t)$ satisfies:
\begin{enumerate}
  \item \textbf{Range:} $G(\mathbf{X}_t) \in [0,1]$.
  \item \textbf{Circuit-Breaker:}
        $G(\mathbf{X}_t) = 0 \Leftrightarrow
         \exists\,d : \mathcal{S}_t^{(d)} \leq \mathcal{S}_{\min}^{(d)}$.
  \item \textbf{Full Deployment:}
        $G(\mathbf{X}_t) = 1 \Leftrightarrow
         \forall\,d : \mathcal{S}_t^{(d)} \geq \mathcal{S}_{\mathrm{full}}^{(d)}$.
  \item \textbf{Monotonicity:} $G$ is non-decreasing in each coordinate.
  \item \textbf{Single-Dimension Sufficiency:}
        One critically discharged state halts all capital deployment,
        regardless of the other six dimensions.
\end{enumerate}
\end{theorem}

\section{The Capital Deployment Policy}

\begin{definition}[State-Gated Capital Deployment Policy]
\label{def:deployment-policy}
The optimal state-gated capital deployment policy is:
\begin{equation}
  \pi^*(\mathbf{X}_t, S_t) \;=\;
  \begin{cases}
    0
      & \exists\,d : \mathcal{S}_t^{(d)} \leq \mathcal{S}_{\min}^{(d)}, \\[6pt]
    K_t^{(\mathrm{inv})} \cdot C_{\mathrm{available}} \cdot G(\mathbf{X}_t)
      & \text{otherwise},
  \end{cases}
  \label{eq:deployment-policy}
\end{equation}
where $C_{\mathrm{available}} = W_t - \ell_{\min}$ is the
deployable capital after reserving the liquidity floor.
\end{definition}

\begin{proposition}[Policy is Well-Defined and Bounded]
\label{prop:policy-bounded}
$\pi^*(\mathbf{X}_t, S_t) \in [0, C_{\mathrm{available}})$
for all $\mathbf{X}_t$ and all $t$.  In particular,
$\pi^* < C_{\mathrm{available}}$ strictly since
$K_t^{(\mathrm{inv})} < 1$ whenever $R_t > 0$.
\end{proposition}

% ==============================================================
\chapter{The Composite Risk-Adjusted Objective}
\label{ch:objective}
% ==============================================================

\section{Objective Function}

\begin{definition}[Composite Risk-Adjusted Objective]
\label{def:objective}
The composite objective function is:
\begin{equation}
  J \;=\; \alpha\,G_{\mathrm{ret}}
         \;-\; \beta\,D_{\mathrm{dd}}
         \;-\; \gamma\,V
         \;+\; \delta\,S_{\mathrm{surv}}
         \;+\; \varepsilon\,R_{\mathrm{risk}}
         \;-\; \zeta\,T_{\mathrm{cost}}
         \;-\; \eta\,C_{\mathrm{conc}}
         \;-\; \theta\,L_{\mathrm{excess}},
  \label{eq:composite-objective}
\end{equation}
where $G_{\mathrm{ret}}$ is total portfolio growth, $D_{\mathrm{dd}}$
is maximum drawdown, $V$ is annualised volatility,
$S_{\mathrm{surv}} = \Prob(W_T \geq \ell_{\min})$,
$R_{\mathrm{risk}}$ is the Calmar ratio, and the remaining terms
penalise excessive turnover, concentration, and leverage.
\end{definition}

\begin{remark}[Equivalence to Negative Posterior MSE]
Maximising $J$ is equivalent to minimising the trace of the
posterior covariance $\mathbf{P}_{t|t}$:
\begin{equation}
  \argmax_\pi J(\pi)
  \;=\;
  \argmin_{\hat{\mathbf{X}}}
  \E\!\left[\|\mathbf{X}_t - \hat{\mathbf{X}}\|^2
      \mid \mathbf{z}_{1:t}\right].
\end{equation}
This formalises the MVUE identification at the core of the platform.
\end{remark}

\section{Layered Circuit-Breaker Architecture}

\begin{table}[H]
\centering
\caption{Four-level circuit-breaker architecture.}
\label{tab:circuit-breakers}
\small
\begin{tabularx}{\textwidth}{@{} l X X @{}}
\toprule
\textbf{Level} & \textbf{Trigger Condition} & \textbf{Response} \\
\midrule
1 Advisory    & Any $\mathcal{S}_t^{(d)} < 0.50$
              & Log warning; reduce sizing 20\%. \\
2 Caution     & Any $\mathcal{S}_t^{(d)} < 0.35$
                or two states $< 0.50$
              & Reduce deployment 40\%; tighten stops. \\
3 Alert       & Any $\mathcal{S}_t^{(d)} < 0.25$
                or $\calP_t < 0.30$
              & Reduce deployment 70\%; hedge all risk. \\
4 Emergency   & Any $\mathcal{S}_t^{(d)} < \mathcal{S}_{\min}^{(d)}$
                or $\calP_t < 0.20$
              & Hard stop; no new positions; audit. \\
\bottomrule
\end{tabularx}
\end{table}

% ==============================================================
\chapter{Portfolio Diversification Mandate}
\label{ch:diversification}
% ==============================================================

\section{Diversification as a First-Class Constraint}

Diversification is an \emph{architectural constraint}, not a
preference.  The platform enforces:

\begin{enumerate}
  \item \textbf{Hard constraints:} maximum concentration per asset,
        sector, currency, and geography.
  \item \textbf{Soft constraints:} target correlation bands,
        volatility-contribution budgets, and factor-exposure limits.
  \item \textbf{Dynamic triggers:} rebalancing proposals generated
        when concentration drifts or correlations spike, before any
        new signal-based trade is considered.
\end{enumerate}

\begin{theorem}[Concentration is Sub-Optimal]
\label{thm:concentration-suboptimal}
Under a strictly positive definite covariance matrix and the
composite objective $J$, any single-asset portfolio
$\bfe_k$ is strictly dominated by a diversified portfolio
$\bfw^* \in \Delta^{n-1}$ unless $\Cov(r_k, r_j) \geq \sigma_k^2$
for all $j \neq k$.
\end{theorem}

\begin{theorem}[CVaR Sub-Additivity]
\label{thm:cvar-subadd}
For any portfolio weights $\bfw \in \Delta^{n-1}$ and
confidence level $\alpha \in (0,1)$:
\begin{equation}
  \CVaR_\alpha\!\left(\sum_i w_i L_i\right)
  \;\leq\;
  \sum_i w_i\,\CVaR_\alpha(L_i),
\end{equation}
with strict inequality whenever the losses are not perfectly
co-monotonic.  Cross-asset diversification always reduces tail risk.
\end{theorem}

\section{Regime-Conditional Allocation}

\begin{table}[H]
\centering
\caption{Regime-conditioned asset-class tilts.}
\label{tab:regime-tilts}
\small
\begin{tabularx}{\textwidth}{@{} l X @{}}
\toprule
\textbf{Regime} & \textbf{Asset-Class Orientation} \\
\midrule
Risk-On      & $\uparrow$ Equities, EM, High-Yield Credit, Commodities, Crypto. \\
Risk-Off     & $\uparrow$ Treasuries, Gold, JPY, CHF, Short-Duration Bonds, Cash. \\
Inflationary & $\uparrow$ TIPS, Commodities, Energy, Real Assets, Commodity FX. \\
Recessionary & $\uparrow$ Long-Duration Bonds, Deflation Structures, Utilities. \\
High-Vol / Crisis & $\uparrow$ Volatility Instruments, Tail Hedges, Defensives, Cash. \\
Stagflation  & $\uparrow$ Real Assets, Gold, Short Equities, Commodity Producers. \\
\bottomrule
\end{tabularx}
\end{table}

% ==============================================================
\chapter{Bayesian Agent Reputation and Recursive Reinvestment}
\label{ch:reputation}
% ==============================================================

\section{Bayesian Reputation Update}

\begin{definition}[Beta Prior for Agent Reliability]
\label{def:beta-prior}
The prior distribution over agent $i$'s reliability in regime
$r$ is $\rho_i^{(r)} \sim \mathrm{Beta}(\alpha_i^{(r)}, \beta_i^{(r)})$.
After observing $n$ episodes with $k$ correct predictions:
\begin{align}
  \alpha_i^{(r)\,\prime} &= \alpha_i^{(r)} + k, \\
  \beta_i^{(r)\,\prime}  &= \beta_i^{(r)} + n - k.
\end{align}
The posterior mean (reputation weight) is:
\begin{equation}
  \hat{\rho}_i^{(r)}
  \;=\;
  \frac{\alpha_i^{(r)} + k}
       {\alpha_i^{(r)} + \beta_i^{(r)} + n}.
  \label{eq:posterior-mean}
\end{equation}
\end{definition}

\begin{theorem}[Geometric Posterior Convergence]
\label{thm:posterior-convergence}
As $n \to \infty$, the posterior mean $\hat{\rho}_i^{(r)} \to
\rho_i^{(r)}$ almost surely, and the posterior variance satisfies
$\mathrm{Var}[\rho_i^{(r)} \mid \mathrm{data}] = O(1/n)$.
Agents whose predictions consistently match reality accumulate
high reputation; those whose predictions are consistently wrong
see their Kalman gain contribution shrink toward zero.
\end{theorem}

\section{Recursive Reinvestment and Compounding}

\begin{theorem}[Monotone Compounding]
\label{thm:monotone-compounding}
Under the composite objective $J$ and the state-gated deployment
policy $\pi^*$, the sequence of composite objective values
$\{J(\pi_t)\}_{t=0}^{T}$ is non-decreasing in expectation:
\begin{equation}
  \E\!\left[J(\pi_{t+1}) \mid \calF_t\right] \;\geq\; J(\pi_t).
\end{equation}
The inequality is strict whenever the portfolio crosses a capital
threshold $\tau_m$, activating a higher-tier allocation strategy.
\end{theorem}

% ==============================================================
\chapter{System Integration and Full Architecture}
\label{ch:integration}
% ==============================================================

\section{The Complete Processing Pipeline}

The full system architecture implements the following recursive
loop at each discrete time step $t$:

\begin{enumerate}
  \item \textbf{Observe:} ingest raw market, macro, and news data
        across all four domains.
  \item \textbf{Estimate States:} seven specialised state agents
        produce $\hat{\mathbf{X}}_{t|t}$ via the Kalman update.
  \item \textbf{Gate:} evaluate $G(\mathbf{X}_t)$; circuit-breakers
        fire if any $\mathcal{S}_t^{(d)} \leq \mathcal{S}_{\min}^{(d)}$.
  \item \textbf{Process Signals:} multi-agent ensemble computes
        $S_t$, $C_t$, $U_t$, $D_t$; Kalman gain $K_t$ is derived.
  \item \textbf{Allocate Capital:}
        $\pi^* = K_t \cdot C_{\mathrm{available}} \cdot G(\mathbf{X}_t)$.
  \item \textbf{Risk Gate:} CVaR, MaxDD, and concentration checks;
        hard stop if any limit is breached.
  \item \textbf{Execute:} orders submitted via Alpaca API.
  \item \textbf{Measure:} realised P\&L vs.\ five baselines;
        composite objective $J$ evaluated.
  \item \textbf{Update Reputations:} Bayesian $w_i$ adaptation;
        measurement noise $\hat{\mathbf{R}}_t$ updated online.
  \item \textbf{Return to Step 1.}
\end{enumerate}

\section{Baseline Comparisons}

The platform is benchmarked against five baselines:

\begin{table}[H]
\centering
\caption{Baselines and the mechanism by which the platform
         strictly dominates each.}
\label{tab:baselines}
\small
\begin{tabularx}{\textwidth}{@{} l X X @{}}
\toprule
\textbf{Baseline} & \textbf{Missing Property}
  & \textbf{Dominated Via} \\
\midrule
Buy-and-Hold
  & No drawdown protection; no CVaR gate.
  & Circuit-breaker prevents catastrophic drawdowns. \\
Equal-Weight Rebalancing
  & Ignores all agent signals.
  & Ensemble signal carries strictly positive information value. \\
Static Allocation
  & No regime adaptation.
  & Bayesian weights detect and respond to regime shifts. \\
Single-Strategy
  & No ensemble diversity.
  & Ensemble variance reduction strictly improves Kalman gain. \\
Greedy Maximiser
  & No risk constraints.
  & Quality dominance: controlled growth beats volatile boom-bust. \\
\bottomrule
\end{tabularx}
\end{table}

\begin{theorem}[Optimality of the Multi-Agent System]
\label{thm:main-optimality}
Under the stated assumptions (positive information value, agent
diversity, positive risk aversion), the multi-agent capital
allocation system $\mathcal{M}$ achieves a strictly higher
expected composite objective than any of the five baselines:
\begin{equation}
  \E[J(\mathcal{M})]
  \;>\;
  \max\!\left(
    \E[J(\pi_{\mathrm{BH}})],\;
    \E[J(\pi_{\mathrm{EW}})],\;
    \E[J(\pi_{\mathrm{SA}})],\;
    \E[J(\pi_{\mathrm{SS}})],\;
    \E[J(\pi_{\mathrm{GM}})]
  \right).
\end{equation}
\end{theorem}

% ==============================================================
\chapter*{Summary of Key Results}
\addcontentsline{toc}{chapter}{Summary of Key Results}
% ==============================================================

\begin{center}
\renewcommand{\arraystretch}{1.35}
\begin{tabularx}{\textwidth}{@{} l X X @{}}
\toprule
\textbf{Result} & \textbf{Key Guarantee} & \textbf{Chapter} \\
\midrule
Theorem~\ref{thm:signal-bounded}
  & $S_t \in [-1,1]$ for all $t$.
  & \ref{ch:agents} \\
Proposition~\ref{prop:gain-properties}
  & $K_t^{(\mathrm{inv})} \in [0,1)$; bounded by confidence and disagreement.
  & \ref{ch:kalman} \\
Theorem~\ref{thm:discharge-dominance}
  & Long-run state charge $< \tfrac{1}{2}$; discharge dominates.
  & \ref{ch:states} \\
Theorem~\ref{thm:product-gating}
  & $G(\mathbf{X}_t) = 0$ iff any dimension breaches its floor.
  & \ref{ch:gating} \\
Proposition~\ref{prop:policy-bounded}
  & $\pi^* \in [0, C_{\mathrm{available}})$; strict sub-capacity.
  & \ref{ch:gating} \\
Theorem~\ref{thm:concentration-suboptimal}
  & Concentration is never variance-minimising in the generic case.
  & \ref{ch:diversification} \\
Theorem~\ref{thm:cvar-subadd}
  & CVaR is sub-additive; diversification strictly reduces tail risk.
  & \ref{ch:diversification} \\
Theorem~\ref{thm:posterior-convergence}
  & Agent reputations converge geometrically to true reliability.
  & \ref{ch:reputation} \\
Theorem~\ref{thm:monotone-compounding}
  & Expected composite objective is non-decreasing under $\pi^*$.
  & \ref{ch:reputation} \\
Theorem~\ref{thm:main-optimality}
  & Multi-agent system strictly dominates all five baselines.
  & \ref{ch:integration} \\
\bottomrule
\end{tabularx}
\end{center}

\medskip
\noindent
\textbf{The deepest result:}

\begin{align}
    \E[J(\mathcal{M})]
    \;&>\;
    \max\!\left(
        \E[J(\pi_{\mathrm{BH}})],\;
        \E[J(\pi_{\mathrm{EW}})],\;
        \E[J(\pi_{\mathrm{SA}})],\;
        \E[J(\pi_{\mathrm{SS}})],\;
        \E[J(\pi_{\mathrm{GM}})]
    \right)
    \nonumber\\[2mm]
    \;&=\;
    \max\!\left(
        \begin{gathered}
        \E[\text{Buy-and-Hold}],\;
        \E[\text{Equal-Weight}],\;
        \E[\text{Static Allocation}],\\
        \E[\text{Single-Strategy}],\;
        \E[\text{Greedy Maximiser}]
        \end{gathered}
    \right).
\end{align}

\noindent
and the finite capital constraint $W_0 = \$100{,}000$ is the
condition that forces the system into the minimum-variance,
survival-first regime --- protecting the principal that makes all
future compounding possible.

% ==============================================================
\chapter*{References}
\addcontentsline{toc}{chapter}{References}
% ==============================================================

\begin{thebibliography}{99}

\bibitem{kalman1960}
Kalman, R.\,E.\ (1960).
A New Approach to Linear Filtering and Prediction Problems.
\textit{Journal of Basic Engineering}, 82(1), 35--45.

\bibitem{markowitz1952}
Markowitz, H.\ (1952).
Portfolio Selection.
\textit{Journal of Finance}, 7(1), 77--91.

\bibitem{rockafellar2000}
Rockafellar, R.\,T.\ \& Uryasev, S.\ (2000).
Optimization of Conditional Value-at-Risk.
\textit{Journal of Risk}, 2(3), 21--41.

\bibitem{barshalom2001}
Bar-Shalom, Y., Li, X.\,R., \& Kirubarajan, T.\ (2001).
\textit{Estimation with Applications to Tracking and Navigation}.
Wiley.

\bibitem{hamilton1989}
Hamilton, J.\,D.\ (1989).
A New Approach to the Economic Analysis of Nonstationary Time Series.
\textit{Econometrica}, 57(2), 357--384.

\bibitem{kelly1956}
Kelly, J.\,L.\ (1956).
A New Interpretation of Information Rate.
\textit{Bell System Technical Journal}, 35(4), 917--926.

\bibitem{plett2004}
Plett, G.\,L.\ (2004).
Extended Kalman Filtering for Battery Management Systems.
\textit{Journal of Power Sources}, 134(2), 252--292.

\bibitem{R8}
\path{finite_investment_architecture_as_economic_financial_fiscal_investment_kalman_filter.md}
--- this repository.
Multi-domain Kalman filter isomorphism; all Kalman notation and
four-domain observation model.

\bibitem{R9}
\path{finite_investment_architecture_states_of_portfolio_investments_securities_finance_markets_fundamentals_sectors.md}
--- this repository.
Seven-state SoC architecture; charge/discharge dynamics;
gating function specification.

\bibitem{R10}
\texttt{high\_level\_finite\_initial\_capital\_master\_capital\_allocation\_proof.tex}
--- this repository.
Master capital allocation proof; ensemble signal; Bayesian
reputation; composite objective; dominance hierarchy.

\bibitem{R11}
\texttt{high\_level\_architecture\_proof.tex} --- this repository.
Operator composition algebra; HMM regime detection; uncertainty
decomposition; CVaR formulation.

\bibitem{R12}
\texttt{high\_level\_supplementary\_diversification\_proof.tex}
--- this repository.
Diversification mandate; CVaR sub-additivity; concentration
constraints; regime-conditional correlation.

\bibitem{R13}
\texttt{high\_level\_kalman\_filter\_states\_capital\_allocation\_proof.tex}
--- this repository.
Full formal proof framework for Kalman filter states and capital
allocation; seven-state gating architecture.

\end{thebibliography}

% ==============================================================
\chapter*{Changelog}
\addcontentsline{toc}{chapter}{Changelog}
% ==============================================================

\begin{center}
\renewcommand{\arraystretch}{1.35}
\begin{tabularx}{\textwidth}{@{} l l l X @{}}
\toprule
\textbf{Version} & \textbf{Date} & \textbf{Author} & \textbf{Description} \\
\midrule
1.0.0 & 2026-08-24 & Hadrian Hu
  & Initial draft.  Established team-mates whitepaper consolidating
    the full multi-agent quantitative-finance platform architecture:
    four-domain Kalman filter isomorphism, seven-state SoC
    architecture, product gating function, Bayesian agent
    reputation, composite risk-adjusted objective, diversification
    mandate, and dominance hierarchy over five baselines.
    Formatted and styled consistent with companion proof monographs.
    \\
\bottomrule
\end{tabularx}
\end{center}

\end{document}
